\documentclass[11pt]{article}
\usepackage[margin=1in]{geometry}
\usepackage{amsmath,amssymb,mathtools}
\usepackage[hidelinks]{hyperref}
\usepackage{microtype}
\setlength{\parindent}{0pt}
\setlength{\parskip}{0.65em}
\newcommand{\HH}{\mathbb H}
\newcommand{\RR}{\mathbb R}
\newcommand{\norm}[1]{\lVert#1\rVert}

\title{A Literature-Implied Upgrade for Quaternion Compressed Sensing}
\author{Run \texttt{fa\_banach\_001} -- lane 1}
\date{26 August 2026}

\begin{document}
\maketitle

\textbf{Status.} Literature-implied answer. This note records a direct
identification with a later block-sparse recovery theorem; it is not claimed
as a new mathematical result.

\section*{Source question}

Bade\'nska and B\l aszczyk prove stable quaternion $\ell_1$ recovery for a
real sensing matrix under $\delta_{2s}<1/3$ \cite[Theorem 4.1 and Corollary
5.1]{BB}. Immediately afterward, on PDF page 12, they write that they
conjecture this requirement is not optimal.

\section*{Supporting theorem}

Gao and Ma prove that mixed $\ell_2/\ell_1$ minimization stably recovers block
sparse signals when the block restricted isometry constant satisfies
\[
  \delta_{2s\mid\mathcal I}<\frac{4}{\sqrt{41}}\approx0.6246,
\]
and in the noiseless, exactly block-$s$-sparse case the recovery is exact
\cite[Theorem 1 and Corollary 1, PDF p.~3]{GM}.

\section*{Identification}

Write each quaternion coordinate as
\[
  z_j=a_j+b_j\mathbf i+c_j\mathbf j+d_j\mathbf k
\]
and let $R:\HH^n\to\RR^{4n}$ collect $(a_j,b_j,c_j,d_j)$ as the $j$th
four-dimensional block. Then
\[
 \norm{z}_{\ell_1(\HH^n)}
 =\sum_{j=1}^n |z_j|
 =\sum_{j=1}^n\norm{Rz[j]}_2
 =\norm{Rz}_{2,1}.
\]
If $\Phi\in\RR^{m\times n}$, realification of $z\mapsto\Phi z$ is, up to a
coordinate permutation,
\[
 A=\Phi\otimes I_4:\RR^{4n}\longrightarrow\RR^{4m}.
\]
It preserves Euclidean data and noise norms. Hence the quaternion program in
\cite{BB} is precisely Gao--Ma's mixed $\ell_2/\ell_1$ program for the block
partition $\mathcal I=(4,\ldots,4)$.

It remains only to compare the two restricted isometry constants. If $u$ is
supported on at most $k$ blocks and $u^{(0)},\ldots,u^{(3)}\in\RR^n$ are its
four component vectors, then each $u^{(r)}$ is $k$-sparse and
\[
 \norm{Au}_2^2=\sum_{r=0}^3\norm{\Phi u^{(r)}}_2^2,
 \qquad
 \norm{u}_2^2=\sum_{r=0}^3\norm{u^{(r)}}_2^2.
\]
Summing the scalar RIP inequalities shows
$\delta_{k\mid\mathcal I}(A)\leq\delta_k(\Phi)$. Conversely, embed any real
$k$-sparse vector in the first component and set the other three components to
zero. This gives the reverse inequality, and therefore
\[
 \delta_{k\mid\mathcal I}(\Phi\otimes I_4)=\delta_k(\Phi).
\]

Applying Gao--Ma with $k=2s$ proves stable quaternion recovery, and exact
recovery in the noiseless sparse case, whenever
\[
 \boxed{\delta_{2s}(\Phi)<4/\sqrt{41}}.
\]
Since $4/\sqrt{41}>1/3$, this fully confirms the source's stated
nonoptimality conjecture in its real-measurement setting.

\section*{Provenance and scope}

This is an agent-identified implication. Gao and Ma do not mention quaternion
signals or the source paper. The threshold $4/\sqrt{41}$ is recorded only as
a certified strict improvement, not as the optimal threshold. Quaternion-valued
sensing matrices are outside the scope.

A bounded search used the exact source title, ``quaternion compressed sensing
RIP'', ``block sparse recovery'', and the $\delta_{2s}$ threshold terminology.
It found Gao--Ma's theorem and later block-RIP refinements, but no paper
explicitly advertising this quaternion specialization.

\begin{thebibliography}{9}
\bibitem{BB}
A.~Bade\'nska and \L.~B\l aszczyk,
\emph{Compressed sensing for real measurements of quaternion signals},
arXiv:1605.07985; J. Franklin Inst. 354 (2017), 5753--5769.

\bibitem{GM}
Y.~Gao and M.~Ma,
\emph{A new bound on the block restricted isometry constant in compressed sensing},
J. Inequal. Appl. 2017, article 174,
\href{https://doi.org/10.1186/s13660-017-1448-2}{doi:10.1186/s13660-017-1448-2}.
\end{thebibliography}

\end{document}

